% -------------------------------------------------------
%  Section 7: Critical assessment of the paper
% -------------------------------------------------------
\section{نقد و ارزیابی انتقادی مقاله}
\label{sec:critique}

\begin{table}[!ht]
    \centering
    \scriptsize
    \caption{دسته‌بندی ایراد‌های استخراج‌شده از متن مقاله}
    \label{tab:critique}
    \begin{tabular}{|p{5.2cm}|p{2.5cm}|p{6.1cm}|}
    \hline
    \textbf{مورد} & \textbf{نوع} & \textbf{شاهد در خود مقاله} \\ \hline
    نبود تکرار و بازه اطمینان برای رتبه‌بندی ۲۰ پیکربندی & خطای روش‌شناختی & جدول ۵؛ پنج ردیف برتر در بازه \num{۰٫۰۲۶۵}--\num{۰٫۰۲۹۲} \\ \hline
    مقایسه دو معماری با چهار ابرپارامتر متفاوت & خطای روش‌شناختی & جدول ۴ \\ \hline
    درهم‌آمیختگی اثر بازنمایی با از دست رفتن پیش‌آموزش لایه ورودی & خطای روش‌شناختی & بخش ۵ (محدودیت‌ها) و جدول ۵ \\ \hline
    ادعای مقاومت در برابر مبهم‌سازی بدون آزمایش & خطای روش‌شناختی & چکیده و بخش ۶ در برابر بخش ۴ \\ \hline
    سنجه واحد و حساس به نامتوازنی کلاس‌ها & خطای روش‌شناختی & جدول ۲ در برابر بخش \num{۴٫۴} \\ \hline
    سیاست کانال‌های غایب مشخص نیست & کاستی گزارش‌دهی & بخش \num{۴٫۲} \\ \hline
    نبود بذر تصادفی، زمان اجرا و سنجه‌های درون‌کلاسی & کاستی گزارش‌دهی & بخش \num{۴٫۴} و جدول ۵ \\ \hline
    نبود مقایسه عددی با کارهای پیشین & کاستی گزارش‌دهی & بخش ۶ در برابر بخش ۲ \\ \hline
    در دسترس نبودن مخزن کد و داده اعلام‌شده & کاستی گزارش‌دهی & پانویس صفحه ۱ \\ \hline
    وابستگی به دیس‌اسمبلر برای استخراج جدول بخش‌ها & محدودیت دامنه & بخش \num{۴٫۱} \\ \hline
    ارزیابی روی یک مجموعه‌داده و تنها قالب \lr{PE} & محدودیت دامنه & بخش \num{۴٫۱} \\ \hline
    \end{tabular}
\end{table}

\subsection{نقاط قوت}
ایده مرکزی ساده، تازه و به‌خوبی انگیزه‌مند است و جای خالی مشخصی را در ادبیات پر می‌کند. سنجه ارزیابی بیرونی با برچسب‌های پنهان (سرور کگل) امکان نشت داده آزمون را از میان می‌برد و ثابت نگه داشتن عرض پایه در \lr{S4} و \lr{S5}، اثر عرض را از اثر قطعه‌بندی جدا می‌کند. بخش محدودیت‌ها صادقانه نوشته شده و روش چنان روشن توصیف شده که بازپیاده‌سازی مستقل از روی متن ممکن بود.

\subsection{خطا‌های روش‌شناختی}

\paragraph{الف. رتبه‌بندی‌ها پشتوانه آماری ندارند.}
۲۰ پیکربندی روی \emph{همان} مجموعه آزمون ارزیابی و سپس بهترین‌ها برجسته شده‌اند. پنج پیکربندی برتر \lr{ResNet50} در بازه \num{۰٫۰۲۶۵} تا \num{۰٫۰۲۹۲} (پهنایی حدود ده درصد نسبی) قرار دارند و هیچ تکرار، انحراف معیار یا بازه اطمینانی گزارش نشده است. پس گزاره «\lr{S3} بهتر از \lr{S1} است» که یکی از دو نتیجه‌گیری اصلی مقاله است، از نوسان تصادفی تفکیک‌پذیر نیست؛ ایراد متوجه \emph{قطعیت لحن} است نه اصل مشاهده.

\paragraph{ب. مقایسه معماری‌ها مخدوش است.}
دو معماری در تعداد دوره، اندازه دسته، نرخ یادگیری اولیه و وجود زمان‌بند تفاوت دارند (جدول \ref{tab:hyper})؛ پس جمله‌هایی از نوع «\lr{S3} روی \lr{ResNet50} به‌مراتب بهتر از \lr{VGG16} عمل می‌کند» دو رویه آموزشی را مقایسه می‌کنند، نه دو معماری را.

\paragraph{پ. اثر بازنمایی با پیش‌آموزش درهم‌آمیخته است.}
در \lr{S1} تا \lr{S3} لایه پیچشی نخست وزن‌های \lr{ImageNet} را دست‌نخورده نگه می‌دارد، ولی در پیکربندی‌های چهار و پنج‌کاناله از نو مقداردهی می‌شود. مقاله این را در بخش محدودیت‌ها می‌پذیرد اما در نتیجه‌گیری از جدول ۵ لحاظ نمی‌کند؛ پس اختلاف \lr{S4} چندکاناله با \lr{S1}--\lr{S3} را نمی‌توان صرفاً به بازنمایی نسبت داد. حتی در حالت سه‌کاناله نیز سه زیرتصویر با معنای کاملاً متفاوت از فیلتر‌هایی عبور می‌کنند که برای همبستگی کانال‌های رنگی آموخته شده‌اند.

\paragraph{ت. ادعای محوری مقاومت، سنجیده نشده است.}
چکیده و نتیجه‌گیری بر «استقلال فضایی» و امید به مقاومت در برابر مبهم‌سازی تأکید می‌کنند، ولی این ادعا سنجیده نمی‌شود؛ در حالی که آزمون لازم پرهزینه نبود: کافی بود ترتیب بخش‌ها در نمونه‌های آزمون جابه‌جا یا داده زائد تزریق شود و افت \lr{S1}--\lr{S3} با \lr{S4} مقایسه گردد. مقاله می‌نویسد «اعتبارسنجی بیشتری لازم است» و این صداقت ارزشمند است، اما چارچوب‌بندی چکیده قوی‌تر از شواهد است. شاهد غیرمستقیمی هم در دست است: چارچوب بعدی همین نویسنده که در بهمن ۱۴۰۴ منتشر شده، آزمون مقاومت در برابر بازآرایی بخش‌ها را صراحتاً فهرست می‌کند.

\paragraph{ث. سنجه با ساختار مسئله همخوان نیست.}
زیان لگاریتمی بر نمونه میانگین می‌گیرد، پس در این مجموعه‌داده نامتوازن، عملکرد روی خانواده‌های کوچک اثری بر عدد نهایی ندارد؛ مدل می‌تواند \lr{Simda} را با \num{۴۲} نمونه به‌کلی از دست بدهد و همچنان زیان لگاریتمی عالی بدهد، در حالی که افزودن \lr{F1} کلان هزینه‌ای نداشت.

\subsection{کاستی‌های گزارش‌دهی}
ستون صحت آموزش برای هر ۲۰ ردیف اشباع است؛ خود مقاله می‌نویسد این سنجه گویا نیست ولی آن را نگه می‌دارد. سیاست برخورد با نمونه‌های فاقد یک بخش مشخص نشده و بذر تصادفی، زمان آموزش و استنتاج و سنجه‌های درون‌کلاسی گزارش نشده‌اند. هیچ عدد مرجعی از کارهای پیشین نیامده در حالی که نتیجه‌گیری ادعای «کارایی قابل مقایسه» دارد، و وعده تحلیل \emph{چرایی} اثر عرض هم برآورده نمی‌شود. جدی‌تر از همه، نشانی مخزن عمومی کد و داده که در پانویس صفحه نخست اعلام شده در دسترس نیست.

\subsection{محدودیت‌های دامنه و جمع‌بندی}
ارزیابی تنها روی یک مجموعه‌داده و قالب \lr{PE} انجام شده و معناشناسی بخش‌ها (به‌ویژه نقش \kd{.rsrc}) ویژه ویندوز است و منتقل نمی‌شود؛ \lr{BIG 2015} هیچ نمونه سالمی ندارد، پس نتایج چیزی درباره نرخ هشدار کاذب نمی‌گویند؛ و وابستگی به دیس‌اسمبلر (\S\ref{sec:dataset}) بخشی از هزینه‌ای را که مصورسازی قرار بود حذف کند بازمی‌گرداند.

جدول \ref{tab:critique} ایراد‌ها را دسته‌بندی می‌کند. ارزیابی کلی ما این است که ایده ارزش پیگیری دارد ولی شواهد، بار ادعا‌ها را تحمل نمی‌کنند و دو پرسش کلیدی بی‌پاسخ می‌ماند: آیا قطعه‌بندی واقعاً مقاومت بیشتری در برابر مبهم‌سازی می‌دهد، و آیا به دامنه‌های دیگر منتقل می‌شود؛ پرسش دوم را در بخش دوم تجربی می‌آزماییم.
